\chapter[Discussão e Conclusão]{Discussão e Conclusão}
\label{cap:conclusao}

\section{Retomada do problema e dos objetivos}

O trabalho atacou a verificação de parentesco facial sob um único protocolo experimental, com o objetivo de comparar arquiteturas representativas da literatura recente e medir o ganho efetivo da combinação proposta neste TCC. O problema é binário, dado um par de imagens decidir se as duas pessoas têm vínculo familiar, e a tarefa é difícil porque parentes próximos podem aparecer com traços distintos enquanto desconhecidos podem ter rostos parecidos. As escolhas de método cobriram cinco direções complementares: baseline sem treinamento com extrator facial pronto, ViT full FT com cross-attention bidirecional, DINOv2 self-supervised com adaptação por LoRA, retrieval-augmented sobre galeria de pares positivos, e a proposta autoral do Modelo 12 com AdaFace parcialmente descongelado, tokens regionais anatômicos, cabeça auxiliar de relação e forward simétrico. A introdução situou o trabalho contra um campo que cresceu entre 2010 e 2025, e os capítulos seguintes detalharam o que cada modelo faz e como foi avaliado.

\section{Discussão dos achados principais}

O resultado central do conjunto experimental é que a combinação de backbone especializado em reconhecimento facial, partial unfreeze do último estágio, comparação regional e simetrização da ordem do par foi a direção mais bem-sucedida observada neste estudo. O Modelo 12 R006 atingiu Test AUC 0,879 e reduziu o gap val→teste para -0,026. O ganho sobre o Modelo 02 R031, com ViT full FT, é de +0,029 em AUC, e o ganho sobre a primeira versão competitiva do próprio Modelo 12 (R002, AUC 0,856) é de +0,022. Essa diferença já supera o piso de ruído de aproximadamente 0,009 AUC estimado no ciclo de reexecuções, embora ainda precise de confirmação por múltiplas seeds para virar afirmação estatística forte.

O baseline B0, AdaFace IR-101 congelado com cosseno, atingiu Test AUC 0,799 sem qualquer treinamento específico para parentesco. O número continua alto o suficiente para mostrar que parte expressiva do sinal de parentesco em FIW já está codificada em embeddings de reconhecimento facial treinados em grandes coleções. O ganho efetivo do Modelo 12 R006 sobre esse extrator facial pronto é de +0,080 AUC. Ao mesmo tempo, B0 ainda produz o maior TAR@FAR=0,001 absoluto do estudo (0,071), acima das variantes treinadas do Modelo 12, o que reforça que operating points de falsa aceitação extremamente baixa exigem análise própria.

O Modelo 02 mostra que o caminho do ViT full FT com cross-attention bidirecional continua competitivo, mas deixa de ser o melhor resultado interno do projeto após o ciclo R005-R010 do Modelo 12. O Modelo 05 com LoRA mantém relevância em eficiência paramétrica, especialmente em regimes de alta precisão, e o Modelo 06 retrieval-augmented ficou abaixo dos paramétricos em AUC absoluto, mas segue como evidência de que memória não-paramétrica pode regularizar classes raras. Sobre os VLMs zero-shot, Codex e Claude ficam em torno de 33\% a 37\% de accuracy em classificação multiclasse, com 0\% nas quatro classes de segundo grau; a tarefa atacada por esses modelos é diferente da verificação binária dos supervisionados, e por isso a comparação é qualitativa.

\section{O que o Modelo 12 realmente demonstra}

Três achados sustentam a contribuição do Modelo 12, e três ressalvas qualificam a leitura. Quanto aos achados, a R006 lidera Test AUC e F1, a R009 lidera TAR@FAR=0,001 entre as variantes treinadas do Modelo 12, e a R010 lidera Average Precision, TAR@FAR=0,01 e TAR@FAR=0,1 dentro da família do Modelo 12. O ganho sobre B0 (+0,080 AUC) e sobre o ViT full FT (+0,029 AUC) indica que a estrutura regional, o partial unfreeze, a cabeça auxiliar de relação e o forward simétrico compõem uma melhoria mensurável. A análise por relação mostra a recuperação mais importante do trabalho: as quatro classes de avós e netos saem da faixa de 36--52\% na R002 para 76--83\% na R006, e a classe gfgs chega a 85{,}7\% na R009.

Quanto às ressalvas, três pontos limitam a força da conclusão. O resultado ainda vem de single-runs em FIW Track-I, sem múltiplas seeds ou k-fold em FIW; o piso de ruído de 0,009 AUC ajuda a interpretar o ciclo, mas não substitui teste estatístico formal. As variantes R006, R009 e R010 otimizam regimes operacionais diferentes, então não há um único ``melhor M12'' para todas as aplicações. Por fim, a contribuição arquitetônica ainda não foi isolada contra baselines B1 e B2 sem componentes regionais, embora o ciclo R005-R010 tenha isolado contribuições internas importantes, especialmente o forward simétrico e a remoção dos vetores globais crus.

\section{Limitações metodológicas}

Esta seção reúne, em uma leitura única, as limitações que a banca poderia identificar, para que o leitor não precise descobrir cada uma por conta própria.

A unidade experimental em FIW é a melhor checkpoint salva em validação dentro de um único run por modelo. KinFaceW-I tem validação cruzada de cinco dobras para o Modelo 02, e o trabalho assume que o split Track-I do RFIW serve como protocolo padrão para FIW. Não há reportagem de múltiplas seeds nem de cross-dataset entre KinFaceW-I e FIW, e a literatura registra perdas significativas em generalização cross-dataset, frequentemente entre 15\% e 30\% de accuracy. O ciclo R003-R010 estimou um piso de ruído prático de aproximadamente 0,009 AUC para alterações de run e sampler, mas isso não substitui variância multi-seed formal. O baseline B0 atende ao §5 do documento de baselines do projeto, mas os baselines B1 (AdaFace global-MLP partial FT) e B2 (AdaFace global-MLP frozen) recomendados pela proposta para isolar a contribuição regional do Modelo 12 ainda não foram implementados nesta versão do trabalho.

A estratégia de invariância geracional por síntese de idade foi planejada como linha transversal aos modelos supervisionados, mas só o Modelo 01 teve o módulo SAM integrado ao pipeline, e mesmo nele apenas parte dos runs operou com o módulo ativo por restrição de hardware (os pesos pré-treinados do SAM têm cerca de 4,5 GB). O Modelo 12 não usa SAM e não participa do esquema diagonal de quatro pares descrito na Seção~\ref{sec:invariancia_geracional} do Capítulo~\ref{cap:Metodologia}. A apresentação correta dessa estratégia neste trabalho é como linha parcialmente executada, não como protocolo uniforme aplicado aos cinco modelos.

As variantes do Modelo 12 evidenciam que o operating point altera fortemente a leitura por classe. A R006 é melhor em AUC agregado e recall, a R009 melhora o regime estrito de baixa falsa aceitação, e a R010 melhora Average Precision e TAR@FAR=0,01. Métricas independentes de threshold (AUC, Average Precision, TAR@FAR) permitem comparação mais direta com os demais modelos, mas accuracy, F1, precision e recall dependem do threshold escolhido em validação e devem ser interpretadas dentro do regime operacional declarado.

A avaliação dos VLMs corresponde a tarefa diferente da verificação binária, classificação multiclasse das onze relações do FIW. Os números reportados (33\% a 37\% de accuracy) servem como leitura qualitativa sobre o estado atual de modelos multimodais generalistas em parentesco facial, não como baseline diretamente comparável aos modelos supervisionados. O experimento de Claude Sonnet ficou restrito a 75 pares por restrição de tempo e custo de chamadas, então o número é piloto e não comparação robusta. Falhas catastróficas em relações de avós e netos (0\% nos VLMs) merecem investigação dedicada com protocolo de amostragem estratificada antes de generalização.

Não foram conduzidas medições de viés demográfico nem testes estatísticos formais para comparar modelos. A literatura aponta diferenças de 5\% a 10\% em taxa de erro entre subgrupos demográficos em modelos de reconhecimento facial, e essa análise não foi reproduzida em parentesco neste trabalho.

\section{Implicações práticas e éticas}

Verificação automática de parentesco facial é uma tecnologia com potencial humanitário e forense relevante, mas com riscos altos se aplicada sem salvaguardas. A introdução referiu aplicações como triagem em rastreamento de crianças desaparecidas, apoio a investigações de tráfico humano e identificação em contextos de conflito armado, e a literatura registra incidentes em que dados genéticos têm sido o instrumento principal de rastreamento.

A leitura responsável do estado atual deste trabalho é que análise facial automática não substitui DNA e não substitui decisão humana qualificada. AUC de 0,879 sobre FIW Track-I, embora seja o maior observado neste estudo, ainda corresponde a taxas de falsos positivos e falsos negativos suficientes para que o uso direto de um modelo deste tipo em decisões irreversíveis seja problemático. A função realista é triagem auxiliar dentro de pipelines com validação humana, particularmente em programas de family tracing conduzidos por organizações como CICV, ACNUR e UNICEF, onde redução do espaço de busca antes de verificação genética pode acelerar identificações sem assumir o papel da prova decisória.

Privacidade biométrica é outra dimensão sensível. A base de imagens deste trabalho é pública e composta por figuras com presença mediática, o que não elimina a obrigação de cuidado com derivações de uso a populações que não consentiram com captura biométrica. Aplicações em contextos forenses devem observar controles institucionais, traceabilidade da decisão e direito à contestação por parte dos indivíduos afetados.

\section{Trabalhos futuros}

A primeira frente de continuidade é a confirmação multi-seed do resultado do Modelo 12 em FIW Track-I, com ao menos três seeds independentes, e a execução de validação cruzada de cinco dobras em FIW para estimar o intervalo de confiança da diferença de AUC para o Modelo 02. A prioridade não é apenas revalidar a R006, mas também verificar se os perfis R009 e R010 se mantêm em baixa falsa aceitação e em mid-FAR quando a amostragem muda.

A segunda frente é a implementação dos baselines B1 e B2 da proposta de baselines do projeto, AdaFace global-MLP com partial FT e com backbone frozen, respectivamente. Eles permitem isolar a contribuição arquitetônica do Modelo 12 contra um modelo de capacidade comparável sem os componentes regionais. Em paralelo, ablações controladas dentro do Modelo 12 (sem cross-attention, sem gate, leave-one-region-out e sem cabeça auxiliar de relação) trariam evidência sobre quais componentes carregam o sinal.

A terceira frente é o tratamento das classes de avós e netos. O gargalo deixou de ser tão severo após R006/R009, mas ainda não está resolvido de forma definitiva: as classes de segundo grau melhoraram muito em accuracy, porém variam entre runs e thresholds, e uma delas ainda regride na R010. Ataques específicos podem incluir amostragem balanceada por relação no batch, retrieval-augmentation localizada apenas para classes raras, ou uma correção real do sampler relation-matched, já que a auditoria do ciclo R004 mostrou que a implementação usada não gerava negativos pareados por relação como pretendido.

A quarta frente é fairness e generalização cross-dataset. Medições de erro condicional a subgrupos demográficos do FIW, e avaliação cruzada entre KinFaceW-I e FIW sob protocolo unificado, são lacunas conhecidas deste trabalho e da literatura recente. Em paralelo, uma avaliação mais cuidadosa dos VLMs em verificação binária pura, com prompts mais bem controlados e amostras maiores, pode esclarecer se a queda observada na adaptação por prompt deste trabalho é específica do Codex ou se reflete um padrão mais amplo da família de modelos multimodais.

\section{Conclusão}

O Modelo 12 obteve o maior Test AUC observado neste estudo com a Run 006 (0,879), superando o melhor modelo anterior, ViT com cross-attention bidirecional FaCoR, por +0,029 AUC. O ganho sobre o baseline sem treinamento (B0, AdaFace cosseno congelado, AUC 0,799) é de +0,080, sustentando a contribuição da combinação entre backbone facial especializado, partial unfreeze do último estágio, comparação regional anatômica, cabeça auxiliar de relação e forward simétrico. A diferença para o ViT full FT é maior do que o piso de ruído observado no ciclo do Modelo 12, mas ainda requer confirmação por múltiplas seeds antes de afirmações categóricas de superioridade estatística.

A leitura conjunta dos resultados sustenta cinco afirmações principais. A especialização do backbone para reconhecimento facial, presente no AdaFace e no esquema regional do Modelo 12, é uma direção produtiva. Atalhos arquitetônicos dependentes da ordem do par podem inflar validação e prejudicar generalização, e o forward simétrico foi a intervenção mais efetiva do projeto para removê-los. Adaptação parameter-efficient via LoRA é competitiva em regimes de alta precisão. Memória não-paramétrica via retrieval regulariza classes raras com perfil per-relação mais uniforme. E modelos multimodais generalistas em zero-shot ainda não substituem treinamento supervisionado específico para parentesco facial.

A contribuição deste trabalho é, portanto, comparativa e metodológica. Mais do que propor uma arquitetura inédita do zero, o estudo aplicou cinco famílias de método distintas sob um único protocolo, mediu o ganho efetivo contra um baseline de referência sem treinamento (B0) e identificou os limites práticos das direções exploradas. O melhor resultado observado deixa de ser apenas incremental sobre o estado da arte interno do projeto: a R006 altera a leitura do Modelo 12 ao fechar boa parte do gap val→teste e recuperar as classes de avós e netos. A agenda concreta de continuação passa por multi-seed, baselines B1/B2, ablações regionais controladas, fairness, correção do sampler relation-matched e tratamento dedicado de regimes de baixa falsa aceitação.
